% 两阶段单纯形法的流程与三种结局：第一阶段判定可行性（正最优值给出Farkas证书），第二阶段给出最优证书或无下界射线。
\begin{tikzpicture}[
  every node/.style={font=\fontsize{8}{10}\selectfont},
  bx/.style={draw=FigureBlue,fill=FigureBlueFill,rounded corners=1.8pt,inner sep=3.5pt,align=center,text width=#1},
  bx/.default=40mm,
  ok/.style={draw=FigureGreen,fill=FigureGreenFill,rounded corners=1.8pt,inner sep=3.5pt,align=center,text width=31mm},
  bad/.style={draw=FigureRed,fill=FigureRedFill,rounded corners=1.8pt,inner sep=3.5pt,align=center,text width=31mm},
  arr/.style={-{Stealth[length=2mm]},draw=FigureStroke,line width=.6pt},
  lab/.style={font=\FigureNoteFont\fontsize{7.5}{9}\selectfont,text=FigureMuted}]
  \node[bx] (s) at (0,0) {标准型$\bm A\bm x=\bm b$，$\bm x\ge\bm 0$\\右端改为非负，加人工变量$\bm a$};
  \node[bx] (p1) at (0,-1.45) {第一阶段：$\min\ \bm 1^\top\bm a$\\初始基＝全部人工变量};
  \node[bad] (inf) at (4.95,-1.45) {原问题不可行\\Farkas证书$\bm z$};
  \node[bx] (clean) at (0,-2.95) {人工变量换出基，\\或删去冗余行};
  \node[bx] (p2) at (0,-4.35) {第二阶段：恢复原目标$\bm c$\\沿边换基};
  \node[ok] (opt) at (-2.35,-5.95) {约化费用全非负\\最优$\bm x$与$\bm y$，间隙为$0$};
  \node[bad] (unb) at (2.35,-5.95) {方向$\bm d\ge\bm 0$且$\bm c^\top\bm d<0$\\无下界射线$\bm x+t\bm d$};
  \draw[arr] (s)--(p1);
  \draw[arr] (p1)--node[lab,above]{最优值$>0$}(inf);
  \draw[arr] (p1)--node[lab,right]{最优值$=0$}(clean);
  \draw[arr] (clean)--(p2);
  \draw[arr] (p2.south)--++(0,-0.25)-|(opt.north);
  \draw[arr] (p2.south)--++(0,-0.25)-|(unb.north);
\end{tikzpicture}
